% =========================================================
% METHODS + RESULTS (updated to match ISDS analysis code)
% Source of truth: make_paper_figures.py, make_equations.py,
%                  TUEP_result occupancy / mLTD summaries
% =========================================================

\section{Methods}

\subsection{Dataset and Study Design}

Two EEG datasets were used in this study for model pretraining and latent-state
evaluation, respectively: the Temple University Hospital EEG Corpus (TUEG) and
the Temple University Epilepsy Corpus (TUEP). Following previous EEG foundation
model studies \cite{zhou2026csbrain,jiang2024large,wang2025cbramod}, EEG
recordings were band-pass filtered between 0.3 and 75~Hz, notch filtered at
60~Hz, mapped to the standard 19-channel international 10--20 montage,
resampled to 200~Hz, and segmented into non-overlapping 30-second windows.
Signal amplitudes were normalized to approximately $\pm100~\mu$V to reduce
inter-recording variability.

TUEG was used for self-supervised pretraining of our lightweight EEG encoder
(USE-FM / USFEEG). The pretraining dataset comprised 2,259,998 EEG segments,
corresponding to approximately 18,800 hours of recordings. The model was
implemented using Python 3.12.12 and PyTorch 2.5.1 with CUDA 12.4 and trained
on four NVIDIA A100 GPUs using Distributed Data Parallel (DDP). Pretraining
used the Adam optimizer with an initial learning rate of $1\times10^{-4}$, a
batch size of 64, and mean squared error (MSE) reconstruction loss. Training
was performed for up to 100 epochs with early stopping (patience = 5),
requiring approximately 100 hours in total.

TUEP was subsequently used for latent-state analysis. Recordings were divided
into epilepsy and non-epilepsy groups, with $n=99$ subjects included in each
group for the primary subject-level occupancy analysis. EEG recordings were
processed into fixed-length segments using a common channel configuration and
preprocessing protocol compatible with both encoders. The same TUEP EEG
segments were independently passed through the frozen LUNA-Base and USE-FM
encoders. No fine-tuning or task-specific adaptation was performed on TUEP for
the dynamical analyses reported here.

Clinical labels were used only to define the epilepsy and non-epilepsy groups.
Because no state-level clinical annotations were available, individual latent
states were treated as data-driven representation states and were not assigned
predefined diagnostic or physiological meanings.


\subsection{EEG Foundation Model Encoders}

To investigate whether epilepsy-related latent-state patterns are reproducible
across independently learned EEG representation spaces, we evaluated two
pretrained EEG foundation model encoders: LUNA-Base and USE-FM (also referred
to as USFEEG), our lightweight EEG foundation model. The two models differ in
architecture, model size, and latent dimensionality, providing distinct
representation spaces for cross-encoder comparison.

USE-FM adopts a CNN--Transformer architecture consisting of two one-dimensional
convolutional layers followed by temporal positional encoding and a two-layer
Transformer encoder. The model was pretrained on TUEG using a self-supervised
signal-reconstruction objective, in which the reconstruction loss minimizes
the difference between the input EEG segment and its reconstructed signal. The
encoder contains 1.46 million parameters and produces a 128-dimensional latent
representation for each EEG segment.

LUNA-Base was selected as an independently pretrained EEG foundation model for
comparison. LUNA employs its pretrained representation architecture to generate
segment-level EEG embeddings and contains approximately 7.10 million
parameters, producing a 256-dimensional latent representation.

\begin{table}[t]
\centering
\caption{Comparison of the EEG foundation model encoders evaluated in this study.}
\label{tab:model_compare}
\setlength{\tabcolsep}{6pt}
\renewcommand{\arraystretch}{1.10}
\begin{tabular}{lccc}
\toprule
Model & Parameters & Backbone & Latent Dim. \\
\midrule
USE-FM (Ours) & \textbf{1.46M} & CNN--Transformer & 128 \\
LUNA-Base & 7.10M & LUNA encoder & 256 \\
\bottomrule
\end{tabular}
\end{table}

For the present study, both pretrained encoders were frozen and used only for
representation extraction. Given an EEG segment $\mathbf{x}$, the two encoders
independently generated latent representations
$\mathbf{h}^{(\mathrm{LUNA})}$ and $\mathbf{h}^{(\mathrm{USE})}$.
No encoder parameters were updated using TUEP labels. This frozen-encoder
design allows group-related latent patterns to be compared without introducing
task-specific adaptation.


\subsection{Latent-State Discovery and Occupancy}

For each encoder $e \in \{\mathrm{LUNA},\mathrm{USE}\}$, frozen segment-level
embeddings were independently partitioned into $k$ latent states using
$k$-means clustering. Clustering characteristics were examined across a range
of $k$ values; $k=4$ was selected for the primary analysis reported here
(consistent with the occupancy and mLTD exports used for figure generation).

Let
\[
z_{s,t}^{(e)} \in \{0,\ldots,k-1\}
\]
denote the latent-state assignment of segment $t$ from subject $s$ under
encoder $e$.

For each subject, the occupancy of state $i$ was defined as the proportion of
EEG segments assigned to that state:
\[
o_{s,i}^{(e)}
=
\frac{1}{T_s}
\sum_{t=1}^{T_s}
\mathbb{I}\left[z_{s,t}^{(e)}=i\right],
\]
where $T_s$ denotes the total number of EEG segments from subject $s$.
Group-level mean occupancy and standard deviation were then calculated
separately for the epilepsy and non-epilepsy groups $G \in \{\mathrm{NonEpi},\mathrm{Epi}\}$:
\[
\bar{o}_{i,G}^{(e)}
=
\frac{1}{|G|}
\sum_{s\in G}
o_{s,i}^{(e)}.
\]

The epilepsy-related occupancy change of state $i$ was defined as
\[
\Delta o_i^{(e)}
=
\bar{o}_{i,\mathrm{Epi}}^{(e)}
-
\bar{o}_{i,\mathrm{NonEpi}}^{(e)}.
\]
Positive values indicate increased state occupancy in epilepsy, whereas
negative values indicate reduced occupancy relative to the non-epilepsy group.


\subsection{Cross-Encoder State Alignment (Hungarian)}

Because $k$-means clustering was performed independently within each
representation space, the resulting state indices are arbitrary. For example,
State~1 identified by LUNA cannot be assumed to correspond directly to State~1
identified by USE-FM. Therefore, state alignment was performed before comparing
group-related occupancy patterns and transition matrices across encoders.

In the current implementation, each latent state $i$ is characterized by the
group-level occupancy signature
\[
f_i^{(e)}
=
\left(
\bar{o}_{i,\mathrm{NonEpi}}^{(e)},
\bar{o}_{i,\mathrm{Epi}}^{(e)}
\right).
\]
Pairwise costs were computed as the \textbf{cosine distance} between LUNA and
USE-FM occupancy signatures:
\[
C_{ij}
=
1-\cos\!\left(
f_i^{(\mathrm{LUNA})},
f_j^{(\mathrm{USE})}
\right).
\]
The Hungarian assignment algorithm
(\texttt{scipy.optimize.linear\_sum\_assignment}) was then used to obtain a
one-to-one matching that minimizes the total cost:
\[
\pi^{*}
=
\arg\min_{\pi}
\sum_{i=0}^{k-1}
C_{i,\pi(i)}.
\]
The resulting optimal permutation $\pi^{*}$ reorders USE-FM states to the LUNA
state order. Aligned USE-FM occupancy changes are
\[
\Delta\tilde{o}_i^{(\mathrm{USE})}
=
\Delta o_{\pi^{*}(i)}^{(\mathrm{USE})}.
\]
Aligned USE-FM transition matrices are similarly reordered by simultaneous
row/column permutation,
\[
\tilde{W}_{G}^{(\mathrm{USE})}
=
P_{\pi^{*}}\,
\bar{W}_{G}^{(\mathrm{USE})}\,
P_{\pi^{*}}^{\mathsf{T}}.
\]

Cross-encoder concordance of epilepsy-related occupancy-change patterns was
quantified using Pearson correlation:
\[
r_{\Delta}
=
\mathrm{corr}
\left(
\Delta\mathbf{o}^{(\mathrm{LUNA})},
\Delta\tilde{\mathbf{o}}^{(\mathrm{USE})}
\right).
\]
This alignment establishes a common state ordering for cross-encoder comparison
but does not imply that matched states have identical physiological or clinical
meanings. Because $\pi^{*}$ is estimated from the same group-level occupancy
signatures that define $\Delta o$, $r_{\Delta}$ should be interpreted as
concordance \emph{after occupancy-informed alignment}, not as an independent
validation that the epilepsy--non-epilepsy contrast was recovered without using
that contrast in the matching criterion.


\subsection{Transition / Causal Dynamics and Structural Summaries}

Transition / causal dynamics were estimated with an mLTD-style group-lasso
procedure, yielding a subject-level weight matrix
$W_s^{(e)}\in\mathbb{R}^{k\times k}$. Group-mean matrices
$\bar{W}_{G}^{(e)}$ were computed for epilepsy and non-epilepsy subjects.
After applying $\pi^{*}$, cross-encoder concordance was quantified by Pearson
correlation of vectorized matrices and by Frobenius distance:
\begin{align*}
r_{W}
&=
\mathrm{corr}\!\left(
\mathrm{vec}(\bar{W}_{G}^{(\mathrm{LUNA})}),
\mathrm{vec}(\tilde{W}_{G}^{(\mathrm{USE})})
\right),
\\
\left\|
\bar{W}_{G}^{(\mathrm{LUNA})}
-
\tilde{W}_{G}^{(\mathrm{USE})}
\right\|_{F}
&=
\sqrt{
\sum_{i,j}
\left(
\bar{W}_{G,ij}^{(\mathrm{LUNA})}
-
\tilde{W}_{G,ij}^{(\mathrm{USE})}
\right)^{2}
}.
\end{align*}
Difference matrices and their concordance were defined as
\begin{align*}
\Delta W^{(e)}
&=
\bar{W}_{\mathrm{Epi}}^{(e)}
-
\bar{W}_{\mathrm{NonEpi}}^{(e)},
\\
r_{\Delta W}
&=
\mathrm{corr}\!\left(
\mathrm{vec}(\Delta W^{(\mathrm{LUNA})}),
\mathrm{vec}(\Delta\tilde{W}^{(\mathrm{USE})})
\right).
\end{align*}
Structural scalars summarized connectivity density. With threshold $\tau\ge 0$
(pipeline default), we computed the number of nonzero edges, group-mean
nonzero-edge counts, mean diagonal weight, and the epilepsy densification
contrast
\[
\Delta\bar{N}_{\mathrm{nz}}^{(e)}
=
\bar{N}_{\mathrm{nz},\mathrm{Epi}}^{(e)}
-
\bar{N}_{\mathrm{nz},\mathrm{NonEpi}}^{(e)}.
\]


\subsection{Cross-Encoder Concordance Protocol and Discrimination}

In summary, the analysis pipeline (i)~discovers $k$-means states independently
per encoder, (ii)~aligns USE-FM states to LUNA by Hungarian assignment on cosine
distances between occupancy signatures
$f_i=(\bar{o}_{\mathrm{NonEpi}},\bar{o}_{\mathrm{Epi}})$,
(iii)~compares $\Delta$~occupancy, $W$, and $\Delta W$ by correlation /
Frobenius distance, (iv)~compares structural scalars (nonzero edges, diagonal
weight), and (v)~reports AUROC for epilepsy vs.\ non-epilepsy discrimination
from dynamical features under each frozen encoder. Because latent states lack
clinical ground truth, concordance of TUEP state-change signatures across
encoders is treated as a reliability check for exploratory clinical dynamics,
not as proof of clinical meaning or deployment readiness.


% =========================================================
% RESULTS
% =========================================================

\section{Results}

\subsection{State Occupancy at $k=4$}

Figure~\ref{fig:occupancy} summarizes subject-level mean occupancy ($\pm$~SD)
at $k=4$ for the epilepsy and non-epilepsy groups under LUNA and USE-FM.
Before state alignment, the two encoders exhibited substantially different
absolute occupancy profiles (e.g., LUNA places substantial mass on States~1
and~3, whereas USE-FM concentrates on States~2--3 with rare State~1),
indicating that the encoders organize their latent representation spaces
differently. Consistent with this observation, absolute occupancy showed
negligible cross-encoder correspondence for both the non-epilepsy
($r=-0.06$) and epilepsy ($r=0.03$) groups.

Within each representation space, Mann--Whitney tests with false discovery
rate (FDR) correction identified significant LUNA group differences for
States~0 and~2, whereas USE-FM state-wise occupancy contrasts did not remain
significant after FDR correction. These results further indicate that
individual latent-state occupancies are sensitive to the learned representation
space and should not be directly interpreted as equivalent states across
encoders.


\subsection{Aligned Occupancy Change and Concordance}

We next examined the relative occupancy change from non-epilepsy to epilepsy,
\[
\Delta o_i
=
\bar{o}_{i,\mathrm{Epi}}
-
\bar{o}_{i,\mathrm{NonEpi}}.
\]
Following cosine-distance Hungarian alignment of USE-FM states to LUNA, the
epilepsy-related occupancy-change vectors showed high cross-encoder concordance
across the four aligned states ($r_{\Delta}=0.96$). Thus, although the two
encoders exhibited substantially different absolute occupancy profiles, their
aligned representations showed a similar direction and relative pattern of
occupancy change between the non-epilepsy and epilepsy groups.

These findings suggest that relative group-associated occupancy patterns may
provide a more stable basis for cross-encoder comparison than absolute
latent-state identities. Because the state alignment itself was based on
group-level occupancy signatures
$f_i=(\bar{o}_{\mathrm{NonEpi}},\bar{o}_{\mathrm{Epi}})$,
however, the observed $r_{\Delta}$ should be interpreted as concordance
following occupancy-informed alignment rather than as an independent validation
of equivalent clinical states.


\subsection{Transition / Causal Matrices and Structural Summaries}

Figure~\ref{fig:W} shows mean causal weight matrices $W$ for both encoders and
groups (USE-FM reordered to the occupancy alignment). Epilepsy matrices were
moderately concordant across encoders ($r=0.67$), whereas non-epilepsy matrices
were weaker ($r=0.30$). Difference matrices
$\Delta W=\mathrm{Epilepsy}-\mathrm{NonEpi}$ showed near-zero edge-wise
correlation ($r_{\Delta W}=-0.07$), indicating that concordance is clearer in
within-group structure and occupancy shifts than in raw difference-edge maps.
Both encoders exhibited denser epilepsy graphs (more nonzero edges; LUNA /
USE-FM densification contrasts $+3.44$ / $+2.78$) and higher mean diagonal
weights in epilepsy.


\subsection{Cross-Encoder Concordance Summary and Group Discrimination}

Table~\ref{tab:concordance} summarizes concordance metrics after state
alignment. The strongest agreements are aligned $\Delta$ occupancy
($r=0.96$), epilepsy $W$ ($r=0.67$), shared denser-epilepsy structure, and
identical AUROC $=0.81$ for epilepsy vs.\ non-epilepsy discrimination from
dynamical features under each encoder. Absolute occupancy profiles and Diff-$W$
edges remain weakly matched, reinforcing that reliability claims should center
on state-change signatures and discriminative utility under limited
state-level ground truth.

\begin{table}[t]
\centering
\caption{Cross-encoder concordance at $k=4$ after Hungarian state alignment
(cosine distance on occupancy signatures).}
\label{tab:concordance}
\setlength{\tabcolsep}{6pt}
\renewcommand{\arraystretch}{1.12}
\begin{tabular}{lc}
\toprule
Metric & LUNA vs.\ USE-FM \\
\midrule
$\Delta$ occupancy correlation (Epi$-$NonEpi) & 0.96 \\
Occupancy correlation (NonEpi / Epi) & $-0.06$ / 0.03 \\
$W$ matrix correlation (NonEpi / Epi) & 0.30 / 0.67 \\
Diff-$W$ correlation & $-0.07$ \\
$W$ Frobenius distance (NonEpi / Epi) & 0.40 / 0.32 \\
$\Delta$ mean nonzero edges (Epi$-$NonEpi) & $+3.44$ / $+2.78$ \\
AUROC (group discrimination from dynamical features) & 0.81 / 0.81 \\
\bottomrule
\end{tabular}
\end{table}

% Optional: if retaining frozen linear-probe tables from the companion SPMB
% evaluation, keep them clearly separated from the dynamical-feature AUROC
% above (0.81), which is the metric exported in the ISDS TUEP_result AUROC
% curves used by make_paper_figures.py.